\section{A sharp universal degree threshold for blowup trees}
\label{sec:mixed-tree-threshold}

We now prove that every tree of maximum degree three gives a stable
blowup when the nonleaf dimensions equal their valences; the leaf
dimensions can be arbitrary. The argument controls the incoming volume
profiles uniformly under both the one-child and two-child integration
operations. Four neighbours with sufficiently large profiles instead
force the central radial foliation to destabilize. This yields a sharp
universal degree threshold, while allowing the stable degree-four and
degree-five arrangements of Section~\ref{sec:mixed-branching}.
A family with two branching vertices gives further stable arrangements.

\begin{theorem}\label{treemix:subcubic}
Let $G$ be a finite tree with at least one edge and maximum degree at
most three. Set $b_v=d_v$ at every nonleaf vertex, and choose any positive
integer $b_v$ at each leaf. Then the mixed blowup $Y(G,b)$ has
anticanonically stable tangent bundle. If $M=|E(G)|$, then
\[
 n=2M+\sum_{v\text{ leaf}}(b_v-1),\qquad \rho=2M+1.
\]
It is smooth Fano and admits no torus-equivariantly locally trivial
bundle presentation with positive-dimensional complete fibre and base.
The normalized degrees of unassigned original divisors satisfy
\begin{equation}\label{treemix:local-bounds}
 A_v\le\frac{25950132}{26414113}<1\quad(d_v=3),\qquad
 A_v\le\frac{24}{25}<1\quad(d_v=2).
\end{equation}
In particular, every saturated tree of maximum degree three gives such
a stable variety with $n=2M$ and $\rho=n+1$.
\end{theorem}

There is no restriction on the distances between branching vertices or
on the relative sizes of the components of $G-v$. The numerical bounds
in \eqref{treemix:local-bounds} concern the indicated individual divisor
degrees; they do not give a uniform slope gap for all proper subspaces
or for leaves of arbitrary dimension.

In the notation of Section~\ref{sec:intrinsic}, the two internal-vertex
bounds imply
\[
 \sum_{e\ni v}J_{v,e}\le\frac{24558189}{26414113}<1
 \quad(d_v=3),\qquad
 \sum_{e\ni v}J_{v,e}\le\frac{22}{25}<1
 \quad(d_v=2).
\]
The proof therefore bounds the sum of the ordered exceptional
intersections, including their dependence on all farther branches.
Theorem~\ref{mixed:tree} is the additional step that promotes these
local inequalities to stability of the entire tangent bundle.

\subsection{Comparison of simplex integrals}

We use two elementary integral constructions together with Efron's
conditional-sum theorem. A function in this subsection is positive and
continuous on $[0,\infty)$ and constant on $[1,\infty)$. For $k\ge1$ set
\[
 V_{p_1,\ldots,p_k}(t)=
 \int_{\substack{x_i\ge0\\\sum_i x_i\le t}}\prod_i p_i(x_i)\,dx,
 \qquad F_{p_1,\ldots,p_k}(t)=V'_{p_1,\ldots,p_k}(t),
\]
and
\[
 B_{p_1,\ldots,p_k}(t)=
 \int_{\substack{x_i\ge0\\\sum_i x_i\le t}}
       (t-\textstyle\sum_i x_i)\prod_i p_i(x_i)\,dx.
\]
The derivative $F$ is the integral on the sum slice in lattice measure,
obtained by eliminating one coordinate with coefficient one.

\begin{lemma}\label{treemix:convolution}
Let $p_i,q_i$ be functions as above, with $q_i$ log-concave. If every
$p_i/q_i$ is nonincreasing, then
\[
 \frac{B_{p_1,\ldots,p_k}(t)}{B_{q_1,\ldots,q_k}(t)}
 \quad\hbox{and}\quad
 \frac{F_{p_1,\ldots,p_k}(t)}{F_{q_1,\ldots,q_k}(t)}
\]
are nonincreasing for $t>0$, and
\begin{equation}\label{treemix:root-comparison}
 \frac{F_{p_1,\ldots,p_k}(t)}{V_{p_1,\ldots,p_k}(t)}
 \le\frac{F_{q_1,\ldots,q_k}(t)}{V_{q_1,\ldots,q_k}(t)}.
\end{equation}
If every $p_i/q_i$ is nondecreasing, both monotonicity assertions and
inequality \eqref{treemix:root-comparison} reverse.
\end{lemma}

\begin{proof}
Take independent variables $X_i$ with densities proportional to
$e^{-x}q_i(x)$ on $[0,\infty)$, and two further independent exponential
variables of rate one. These densities are log-concave. The bounded
positive function $R(x)=\prod_i p_i(x_i)/q_i(x_i)$ has the asserted
coordinatewise monotonicity and is independent of the two extra
coordinates. Extend it to real arguments by their nonnegative parts.
Efron's conditional-sum theorem
\cite[Theorems~3.2--3.3]{saumardwellner2018} makes its conditional
expectation monotone in the total sum. The nonincreasing version follows
by applying the theorem to $-R$.

On a slice of total $t$, the common exponential factor cancels. Integrating
the two additional coordinates gives the factor $t-\sum_i x_i$, so the
conditional expectation is exactly $B_p(t)/B_q(t)$. Omitting these two
variables instead gives $h(t)=F_p(t)/F_q(t)$; for $k=1$ this is just the
assumed ratio order. Rescaling to a fixed compact simplex proves continuity
of all slice ratios for $t>0$. Their positive denominators ensure that
the monotone versions supplied by the theorem agree with these ratios
everywhere. In the nonincreasing case,
\[
 V_p(t)=\int_0^t F_q(s)h(s)\,ds\ge h(t)V_q(t),\qquad
 F_p(t)=F_q(t)h(t).
\]
Division proves \eqref{treemix:root-comparison}. The same calculation
with the reversed integral inequality proves the other case.
\end{proof}

The following coefficient identity specifies the finite rational
calculations used below. For a function $p$ as above define the formal
power series
\begin{equation}\label{treemix:Q}
 Q_p(z)=p(1)+z\sum_{j\ge0}\frac{z^j}{j!}
                   \int_0^1(1-a)^jp(a)\,da.
\end{equation}
Only finitely many coefficients are needed in each application.

\begin{lemma}\label{treemix:coefficients}
For $k\ge1$, $r\ge0$, and an integer $s\ge0$,
\begin{equation}\label{treemix:split}
 \int_{\substack{a_i\ge0\\\sum_i a_i\le k+r}}
 \frac{(k+r-\sum_i a_i)^s}{s!}\prod_i p_i(a_i)\,da
 =[z^{k+s}]e^{rz}\prod_i Q_{p_i}(z).
\end{equation}
For $r>0$, the facet volume is
\[
 F_{p_1,\ldots,p_k}(k+r)=[z^{k-1}]e^{rz}\prod_i Q_{p_i}(z).
\]
\end{lemma}

\begin{proof}
Partition the domain according to the set $S$ of coordinates in $[0,1]$.
Subtract one from each of the remaining $k-|S|$ coordinates. Their
integration leaves
\[
 \frac{(r+\sum_{i\in S}(1-a_i))^{k-|S|+s}}{(k-|S|+s)!}
 \prod_{i\notin S}p_i(1)
\]
times the product of the short-coordinate weights. Express this power
as an exponential coefficient, integrate the short coordinates over
$[0,1]^S$, and sum over $S$. The result is
\eqref{treemix:split}. Differentiate the case $s=0$ with respect to $r$
to obtain the facet formula.
\end{proof}

\subsection{The proof for maximum degree three}

Fix an oriented edge towards a parent vertex. With the parent's
incidence coordinate equal to $a$, integrate over the component on
the other side of the edge and denote its volume by $p(a)$.
These incoming functions are positive, continuous, and constant for
$a\ge1$. Indeed, the only inequality involving the fixed coordinate is
$a+x\ge1$; above one it is redundant. Positivity follows by choosing
all assigned coordinates slightly greater than one and all remaining
coordinates small and positive. The inequalities $b_v\ge d_v$ leave
room in every vertex simplex. Continuity follows by dominated convergence
inside the fixed compact product of these simplexes. Since removing a
vertex separates the tree into disjoint components, the incoming
functions multiply when its coordinates are fixed.

For a leaf of dimension $b$, the incoming function is
\[
 \phi_b(a)=\frac{(b+\min(a,1))^b}{b!}.
\]
At a nonleaf of degree two or three, integration of the coordinate
assigned to the parent gives, for $0\le a\le1$,
\begin{align}
 T(p)(a)&=B_p(2+a),\label{treemix:T}\\
 \mathcal G(p,q)(a)&=B_{p,q}(3+a),\label{treemix:G}
\end{align}
followed by constant continuation. There are respectively one and two
child coordinates; the remaining interval has length
$2+a-x$ or $3+a-x-y$.

\begin{proof}[Proof of Theorem~\ref{treemix:subcubic}]
Put $P(a)=(2+\min(a,1))^3$, and write $p\preceq P$ if $p/P$ is
nonincreasing. The reference $P$ is an auxiliary function, not the
volume of an additional projective-space factor. It is log-concave:
its logarithmic derivative is $3/(2+a)$ below one and drops to zero
above one. Every leaf satisfies $\phi_b\preceq P$, since
\[
 \frac{d}{da}\log\frac{\phi_b(a)}{P(a)}
 =\frac b{b+a}-\frac3{2+a}\le1-\frac3{2+a}\le0
 \qquad(0<a<1).
\]
This includes every positive leaf dimension.

We next prove closure under \eqref{treemix:T}--\eqref{treemix:G}.
Direct integration, or Lemma~\ref{treemix:coefficients}, gives
\begin{align*}
 T(P)(a)&=\frac{363}{10}+\frac{173}{4}a+\frac{27}{2}a^2,\\
 \mathcal G(P,P)(a)&=\frac{119659}{80}+\frac{148781}{80}a
                  +\frac{3213}{4}a^2+\frac{243}{2}a^3.
\end{align*}
For a polynomial $g$ of degree at most three,
\[
 \left(\frac{g(a)}{(2+a)^3}\right)'
 =\frac{(2+a)g'(a)-3g(a)}{(2+a)^4}.
\]
The numerator coefficients are as follows; the cubic coefficient is zero.
\begin{center}
\begin{adjustbox}{max width=\textwidth,center}
\begin{tabular}{lrrr}
\toprule
$g$ & $1$ & $a$ & $a^2$\\
\midrule
$T(P)$ & $-112/5$ & $-65/2$ & $-27/2$\\
$\mathcal G(P,P)$ & $-12283/16$ & $-20261/40$ & $-297/4$\\
\bottomrule
\end{tabular}
\end{adjustbox}
\end{center}
All are negative. Hence $T(P)\preceq P$ and
$\mathcal G(P,P)\preceq P$. By Lemma~\ref{treemix:convolution},
$T(p)/T(P)$ and $\mathcal G(p,q)/\mathcal G(P,P)$ are nonincreasing
whenever $p,q\preceq P$. Multiplying positive nonincreasing ratios
proves $T(p),\mathcal G(p,q)\preceq P$. Induction on the number of
vertices in each oriented component now gives $p\preceq P$ for every
incoming function. The two child inputs need not be equal.

At a degree-three vertex, the simplex capacity is four. For three
reference inputs $P$, the volume and facet volume are
\[
 V_{P,P,P}(4)=\frac{26414113}{320},\qquad
 F_{P,P,P}(4)=\frac{6487533}{80}.
\]
All arithmetic in these values and the two preceding reference
polynomials follows from
\[
 Q_P(z)=27+\frac{65}{4}z+\frac{131}{20}z^2
                   +\frac{233}{120}z^3+O(z^4)
\]
and Lemma~\ref{treemix:coefficients}: the root values are the coefficients
of $z^3$ and $z^2$ in $e^zQ_P(z)^3$, while the outgoing polynomials
are $[z^2]e^{(1+a)z}Q_P(z)$ and
$[z^3]e^{(1+a)z}Q_P(z)^2$.
Inequality~\eqref{treemix:root-comparison} therefore gives
$A_v\le25950132/26414113$, whose deficit from one is
$463981/26414113>0$.

At a degree-two vertex, each incoming function satisfies
\[
 \frac{p(a)}{p(1)}\ge\frac{P(a)}{P(1)}\ge a
 \quad(0\le a\le1),\qquad
 (2+a)^3-27a=(1-a)^2(a+8)\ge0.
\]
Thus both inputs satisfy precisely \eqref{mixed:incoming-bound}.
The two-input calculation
\eqref{mixed:path-volume}--\eqref{mixed:path-facet}, which requires only
those inequalities and the constant continuation, gives $24V-25F\ge0$.
Hence $A_v\le24/25$. At every leaf, Lemma~\ref{mixed:flux} gives
$A_v<1$, also when its dimension exceeds one. All vertices satisfy
Theorem~\ref{mixed:tree}, which proves stability against all proper
subspaces and hence against all subsheaves. The geometric assertions
and counts follow from Proposition~\ref{mixed:geometry} and
$\sum_vd_v=2M$.
\end{proof}

\subsection{A local obstruction at degree four}

\begin{theorem}\label{treemix:degree-four}
Let $Y(G,b)$ be a Fano mixed blowup indexed by a finite tree, with
$b_u\ge d_u$ at every vertex. Suppose that a vertex $v$ and its four
neighbors all have degree four and factor dimension four. Then
\begin{equation}\label{treemix:four-bound}
 A_v\ge\frac{40090063176909393024}{39498346443818421703}>1.
\end{equation}
The unassigned ray at $v$ defines a strict rank-one destabilizer.
No additional condition is imposed on the branches beyond these neighbors.
\end{theorem}

\begin{proof}
Let $\ell(a)=1+\min(a,1)$. Every incoming branch function $p$ in a Fano
tree satisfies that $p/\ell$ is nondecreasing. To see this, let $b$ be
the dimension of its first vertex, $x$ the coordinate assigned to the
parent, and $z$ its other $b-1$ coordinates. After integration over more
distant vertices, their volume $\Psi(z)$ is coordinatewise nondecreasing,
since increasing a coordinate only relaxes mixed inequalities. It is
independent of $x$. For $0\le a\le1$, integration of $x$ from $1-a$ to
$b+1-\sum z_j$ and rescaling give
\[
 p(a)=(b+a)^b\int_{\substack{y\ge0\\\sum_j y_j\le1}}
                   (1-\textstyle\sum_j y_j)\Psi((b+a)y)\,dy.
\]
The integral factor is nondecreasing; the same formula includes $b=1$
as a zero-dimensional integral. Moreover,
\[
 \frac{d}{da}\log\frac{(b+a)^b}{1+a}
 =\frac{a(b-1)}{(b+a)(1+a)}\ge0.
\]
Both functions are constant above one, which proves the claim.

At a degree-four neighbor $w$ of $v$, the incoming function is
$g_w(a)=B_{p_1,p_2,p_3}(4+a)$ on $[0,1]$, where $p_i$ are its three
child inputs. Apply the increasing case of
Lemma~\ref{treemix:convolution} with reference inputs $\ell$.
This function is log-concave, and each $p_i/\ell$ is nondecreasing.
It follows that $g_w/g_4$ is nondecreasing, where
\begin{equation}\label{treemix:g-four}
 g_4(a)=B_{\ell,\ell,\ell}(4+a)
 =\frac{5217+5965a+2610a^2+520a^3+40a^4}{120}
 \quad(0\le a\le1),
\end{equation}
with constant continuation.

For a further application of the lemma, we check log-concavity of $g_4$.
Writing $H=120g_4$ on $[0,1]$, direct differentiation gives
\begin{align*}
 (H')^2-HH''={}&8348485+14860260a+11120040a^2\\
 &+4474400a^3+1020000a^4+124800a^5+6400a^6.
\end{align*}
Its coefficients are positive. Also $H'>0$, so the logarithmic
derivative drops to zero at the constant continuation. Thus $g_4$ is
log-concave on the whole half-line.

The four components at $v$ supply four functions whose ratios to $g_4$
are nondecreasing. Their product is the outside volume at fixed central
coordinates. The increasing case of
\eqref{treemix:root-comparison} compares $A_v$ with the ratio for four
copies of $g_4$ at capacity five. The reference values are
\[
 V_0(5)=\frac{39498346443818421703}{23224320000},\qquad
 F_0(5)=\frac{104401206189868211}{60480000}.
\]
For their exact evaluation, formula~\eqref{treemix:Q} gives
\[
 Q_{g_4}(z)=\frac{598}{5}+\frac{3683}{48}z+\frac{513}{16}z^2
       +\frac{195899}{20160}z^3+\frac{11549}{5040}z^4+O(z^5).
\]
Lemma~\ref{treemix:coefficients} identifies $V_0(5),F_0(5)$ with
$[z^4]e^zQ_{g_4}(z)^4$ and $[z^3]e^zQ_{g_4}(z)^4$.
Their ratio proves \eqref{treemix:four-bound}; its excess over one is
$591716733090971321/39498346443818421703$.
The unassigned ray is the negative sum of four independent basis
vectors and is the only ray on its line: other original rays lie on
different lines or in different vertex summands, and exceptional rays
have nonzero components in two summands. That line has normalized
weight exactly $A_v>1$, proving strict instability.
\end{proof}

\begin{corollary}\label{treemix:threshold}
Three is the largest integer $D$ for which every saturated mixed-blowup
tree of maximum degree at most $D$ has anticanonically stable tangent
bundle. A counterexample at degree four has dimension $32$ and Picard
number $33$.
\end{corollary}

\begin{proof}
Theorem~\ref{treemix:subcubic} proves the assertion for $D=3$.
For $D=4$, take one degree-four vertex joined to four degree-four
vertices, each with three additional pendant leaves, and saturate every
vertex. There are sixteen edges, five dimension-four factors, and twelve
unit leaves, so $n=32$ and $\rho=33$. All comparisons in
Theorem~\ref{treemix:degree-four} are equalities in this case. Thus
\eqref{treemix:four-bound} gives an unstable example of maximum degree
four, which also excludes every $D\ge4$.
\end{proof}

\subsection{Two branching vertices with arbitrary separation}

The sharp universal threshold does not preclude stable trees containing
degree-four or degree-five vertices. The following smaller family allows
an arbitrary distance between two such vertices, while keeping their
pendant edges of length one.

\begin{proposition}\label{treemix:two-centres}
Let $d,e\in\{3,4,5\}$. Join vertices of degrees $d,e$ by a path of
$L\ge1$ edges, and attach respectively $d-1,e-1$ unit-length pendant
edges. Saturate every vertex. The resulting mixed blowup is smooth Fano,
has anticanonically stable tangent bundle, and admits no nontrivial
torus-equivariantly locally trivial bundle presentation. Writing
$M=d+e+L-2$, its dimension and Picard number are $n=2M$ and $\rho=2M+1$.
\end{proposition}

The proof is in Appendix~\ref{sec:blowup-comparisons}. It bounds the
degree-two vertices uniformly and proves positivity at the two branching
vertices by a linear functional on their incoming slice polynomials.


The fixed pendant lengths in Proposition~\ref{treemix:two-centres} are
essential to its statement. Corollary~\ref{branch:degree-five} prevents
an extension allowing all pendant arms at a degree-five centre to have
arbitrary lengths. The degree-four obstruction, in turn, concerns trees;
its factorization argument does not establish the same assertion for
simple graphs with cycles.
